\hypertarget{appendix-7107009c505b}{}\label{referencestart:appendix-7107009c505b}
\section{Local rigidity and mesh orientation fields}\label{appendix-7107009c505b-local-rigidity-and-mesh-orientation-fields}

\textbf{Curriculum depth:} Developed · \textbf{First seen:} University courses · \textbf{Latest development:} University courses

\begin{allmightroledefinition}

\subsection{Definition and mechanism}\label{appendix-7107009c505b-definition-and-mechanism}

As-Rigid-As-Possible (ARAP) deformation seeks vertex positions that change as little as possible from a preferred local rigid transform. For each vertex neighborhood, fit a rotation to the deformed edge vectors, then minimize residuals between transformed rest edges and current edges. Alternating between local rotations and a global position solve makes a nonlinear shape-editing problem tractable, but the result depends on constraints, initialization, mesh quality, and rotation fitting. Instant Meshes uses direction or orientation fields over a surface to guide a quad-dominant remeshing; field singularities encode where a globally consistent direction cannot exist. A mesh's half-edge connectivity, boundary structure, and orientation determine valid neighborhoods and field transport. Poor triangles, seams, boundary conditions, or inconsistent orientation can create artifacts. ARAP is a local preservation objective, not exact rigidity of the whole shape; a direction field guides topology but does not uniquely determine a good mesh without sizing and feature constraints.

\end{allmightroledefinition}

\subsection{Prerequisites}\label{appendix-7107009c505b-prerequisites}

\begin{itemize}
\tightlist
\item
  \href{../04-Light-Transport-and-Image-Synthesis/main.pdf\#nameddest=appendix-1aea26b2307f}{Half-edge mesh connectivity (Light Transport and Image Synthesis, Reference notes: Half-edge mesh connectivity)}
\item
  \href{../04-Light-Transport-and-Image-Synthesis/main.pdf\#nameddest=appendix-fab348609369}{Manifold polygon meshes and boundary structure (Light Transport and Image Synthesis, Reference notes: Manifold polygon meshes and boundary structure)}
\item
  \href{../04-Light-Transport-and-Image-Synthesis/main.pdf\#nameddest=appendix-0234ab7f86b1}{Subdivision surface masks (Light Transport and Image Synthesis, Reference notes: Subdivision surface masks)}
\item
  \href{../04-Light-Transport-and-Image-Synthesis/main.pdf\#nameddest=appendix-29f4bea4c2bb}{Constrained optimization and inverse kinematics (Light Transport and Image Synthesis, Reference notes: Constrained optimization and inverse kinematics)}
\item
  \href{../07-Vision-XR-and-Playable-Worlds/main.pdf\#nameddest=appendix-2391db8052b1}{Cross products and oriented area (Vision, XR, and Playable Worlds, Reference notes: Cross products and oriented area)}
\end{itemize}

\subsection{Used by}\label{appendix-7107009c505b-used-by}

\begin{itemize}
\tightlist
\item
  \hyperref[03-visual-computing-i-ch002-meshes-connectivity-and-geometry-processing]{Meshes, connectivity, and geometry processing}
\end{itemize}

\subsection{Full teaching explanation}\label{appendix-7107009c505b-full-teaching-explanation}

\begin{itemize}
\tightlist
\item
  \hyperref[03-visual-computing-i-ch002-meshes-connectivity-and-geometry-processing]{Meshes, connectivity, and geometry processing}
\end{itemize}

\subsection{Recall and next depth}\label{appendix-7107009c505b-recall-and-next-depth}

\begin{itemize}
\tightlist
\item
  Describe one local rotation/global solve step of ARAP and the role of a field singularity in remeshing.
\item
  Why can low local deformation energy fail to ensure a globally rigid shape?
\end{itemize}

\subsection{Source and navigation}\label{appendix-7107009c505b-source-and-navigation}

Teaching narrative: \hyperref[03-visual-computing-i-ch002-meshes-connectivity-and-geometry-processing]{open the main explanation}. Browse the \hyperref[appendix-a6d556a4ef50-university-graphics-vision-and-photography---independent-series]{University graphics and vision index}. Additional chapter connections: \hyperref[03-visual-computing-i-ch002-meshes-connectivity-and-geometry-processing]{Meshes, connectivity, and geometry processing}; \hyperref[03-visual-computing-i-ch002-surfaces-patches-topology-and-parameterization]{Surfaces, patches, topology, and parameterization}; \href{../02-Mathematics-for-Visual-Computing/main.pdf\#nameddest=01-foundations-i-ch020-put-the-algorithms-in-a-graphics-context}{Put the algorithms in a graphics context (Mathematics for Visual Computing, Ch. 4)}.

Caption/source records: \href{../07-Vision-XR-and-Playable-Worlds/main.pdf\#nameddest=04-visual-computing-ii-ch036-p3-16---subdivision-surfaces-part-1}{P3-16 (Vision, XR, and Playable Worlds, Ch. 5)}, \href{../07-Vision-XR-and-Playable-Worlds/main.pdf\#nameddest=04-visual-computing-ii-ch036-p3-17---subdivision-surfaces-part-2}{P3-17 (Vision, XR, and Playable Worlds, Ch. 5)}, \href{../07-Vision-XR-and-Playable-Worlds/main.pdf\#nameddest=04-visual-computing-ii-ch036-sc-15---stanford-cs231n-deep-learning-for-computer-vision-spring-2025-lecture-15-3d-vision}{SC-15 (Vision, XR, and Playable Worlds, Ch. 5)}.
